\documentclass[12pt]{article}
\usepackage{samstyle} 

\newcommand{\NuevaPagina}[1]{\newpage
\mbox{
\begin{minipage}{3cm}
    \includegraphics[width=3cm]{logo}
\end{minipage}
\begin{minipage}{14.7cm}
  \begin{flushright}
\hfill P�gina #1 \hfill Primer parcial-ST0240 \hfill Marzo
\hfill C�digo: \rule{4cm}{0.3pt}
  \end{flushright}
\end{minipage}}\vspace{3pt}

\rule{18cm}{1.5pt}}

%%% COMIENZO DE DOCUMENTO %%%
\begin{document}
\pagestyle{empty}
%%% Encabezado
\mbox{
    \begin{minipage}{5cm}
        \includegraphics[width=5cm]{logo}
    \end{minipage}
    \begin{minipage}{13cm}
    \begin{center}
        \textsl{
        \textbf{\Large Escuela de Ingenier�as}\\
        \textbf{\large Programaci�n de Computadores - ST0240}\\
        \textbf{\large Pr�ctica 2 (Valor 20\%)}}
    \end{center}
    \end{minipage}
}
\begin{flushright}
  \begin{minipage}{4.5cm}
    Nota: \hrulefill
  \end{minipage}
\end{flushright}\vspace{-2mm}

% Nombre: \rule{9cm}{0.3pt} \; C�digo: \hrulefill\vspace{6pt}

Profesor: \underbar{\hspace{1cm} Sergio Monsalve
\hspace{3.1cm}} \; \hfill Grupo: \underbar{\hspace{1cm} 063 } \;
Mayo \underbar{\hspace{1cm}} de 2016 \rule{18cm}{1.5pt}
%%% Fin del encabezado

La siguiente es la descripci�n de la Practica 2 del curso de Programaci�n de computadores, esta debe ser entregada m�ximo el d�a domingo 15 de Mayo a las 11:59pm a trav�s de la plataforma de EAFIT Interactiva.\\

Suponga que usted es el profesor(a) de una Materia universitara, tal materia tiene \emph{N} cantidad de actividades evaluativas, que se encuentran indexadas en un archivo dado (\emph{indice.txt}).

\section*{Entrada}

Como entradas usted tiene un archivo <<\emph{indice.txt}>> que contiene:
 
\begin{itemize}

\item  la cantidad de archivos
\item  la cantidad de estudiantes
\end{itemize}

 N veces (1 por cada archivo):

\begin{itemize}
\item El nombre de cada archivo
\item Porcentaje actividad evaluativa
\item La cantidad de estudiantes que presetaron esa actividad evaluativa
\item Que estudiantes no presentaron tal evaluaci�n
\item La cantidad de puntos que tenia cada actividad y los porcentajes de cada punto. 
\end{itemize}

\vdots

Adicionalmente, un archivo por cada evaluaci�n donde esta el identificador de cada estudiante y las notas por cada punto. 

Cantidad de Archivos: N\\
Estudiantes totales: M\\[0.5cm]

Parcial1.csv\\
Porcentaje de la evaluacion: [$valor$ $entre$ $1$ y $100$]\\
Presentaron: $v$ $estudiantes$, donde $v \leq M$\\
Total Puntos: $P$\\
Estudiantes que no presentaron: [$valor \leq M$]\\
porcentaje de los puntos: [$P_1, P_2,...,P_N$]

\newpage

\section*{Salida}

Se requiere que usted elabore un programa que cumpla con los siguientes requisitos: 

\begin{enumerate}
\item \textbf{(Valor 25\%)} Manipulaci�n datos:
\end{enumerate}

\begin{itemize}
  \item Leer los archivos a partir del archivo indice, y obtener todos los datos para procesar
  \item Realizar un arreglo con las notas de cada uno de los estudiantes
  \item Realizar un arreglo por estudiante con las notas de cada evaluaci�n
\end{itemize}
 \vspace{0.5cm}

\begin{enumerate}[resume] 
\item \textbf{(Valor 25\%)}Procesamiento
\end{enumerate}

\begin{itemize}
  \item Encuentre el Promedio de las notas de cada estudiante.
  \item Encuentre la desviaci�n t�pica para cada estudiante.
  \item Encuentre la Nota promedio para cada evaluaci�n
  \item Encuentre la desviaci�n t�pica por evaluaci�n.
\end{itemize}
\vspace{0.5cm}

\begin{enumerate}[resume]
\item \textbf{(Valor 25\%)} Visualizaci�n
\end{enumerate}


\begin{itemize}
  \item Que el programa grafique las notas de cada uno de los estudiantes y los guarde en un archivo con todas las graficas
  \item Graficar el progreso de todos los estudiantes (Gr�fica de lineas con linea por cada estudiante) guardarlo como una imagen aparte
\end{itemize}
\vspace{0.5cm}

\begin{enumerate}[resume]
\item \textbf{(Valor 25\%)} Interpretaci�n y Reportes
\end{enumerate}


\begin{itemize}
  \item Determinar cu�les estudiantes han mejorado su rendimiento y mostrarlo mediante una gr�fica (solo estudiantes que siempre mejoraron su nota)
  \item Exportar las notas tanto de cada evaluacion como finales en un archivo con formato CSV ordenados de mejor a peor nota final. 
\end{itemize}

\end{document}
